\section{Student's $t$ Distribution}
\label{sec:distribucion-t}

The $t$-distribution arises when we standardize a sample mean using the sample standard deviation (which is a random variable) instead of the population standard deviation.

\begin{definicion}[Student's $t$-Distribution]
Let $Z \sim N(0,1)$ and $\chi^2_\nu$ be a chi-square variable with $\nu$ degrees of freedom, independent of each other. Then
\begin{align}
 \label{eq:3.2.12}
 T = \frac{Z}{\sqrt{\chi^2_\nu / \nu}}
\end{align}
follows a \emph{Student's $t$-distribution} with $\nu$ degrees of freedom. We write $T \sim t_\nu$.
\end{definicion}

{Properties}
\begin{itemize}
 \item The probability density function of $T$ is
 \begin{align}
  \label{eq:3.2.13}
  f(t) = \frac{\Gamma\!\left(\frac{\nu+1}{2}\right)}{\sqrt{\nu\pi}\,\Gamma\!\left(\frac{\nu}{2}\right)}
  \left(1 + \frac{t^2}{\nu}\right)^{-(\nu+1)/2}, \quad t \in \mathbb{R}.
 \end{align}
 \item $E(T) = 0$ for $\nu > 1$.
 \item $\text{Var}(T) = \nu/(\nu - 2)$ for $\nu > 2$.
 \item As $\nu \to \infty$, $T \xrightarrow{d} N(0,1)$. For small $\nu$, $T$ has heavier tails than the normal distribution.
\end{itemize}

\begin{observacion}
The $t$-distribution arises when estimating the mean $\mu$ with unknown $\sigma$: if $X_1, \ldots, X_n \sim N(\mu, \sigma^2)$ iid and $S^2$ is the sample variance, then
\begin{align}
 \label{eq:3.2.14}
 \frac{\bar{X} - \mu}{S/\sqrt{n}} \sim t_{n-1}.
\end{align}
This is the \emph{$t$-statistic} that gives its name to the $t$-test (Section 3.6). This is precisely the decomposition $T = Z/\sqrt{\chi^2_{n-1}/(n-1)}$ obtained from Fisher's theorem in the previous section.
\end{observacion}

\subsubsection{Confidence intervals with unknown $\sigma$: the small-sample case}

When $\sigma$ is unknown (the most common practical case), the confidence interval for $\mu$ must use quantiles of the $t$-distribution instead of the standard Normal.

\begin{teorema}[Confidence interval for $\mu$ with unknown $\sigma$]
 \label{eq:5.4.1}
Let $X_1, \ldots, X_n \sim N(\mu, \sigma^2)$ iid with unknown $\sigma^2$. A $(1-\alpha)\times 100\%$ confidence interval for $\mu$ is
\begin{align}
 \bar{X} \pm t_{n-1,\,\alpha/2} \cdot \frac{S}{\sqrt{n}},
\end{align}
where $t_{n-1,\,\alpha/2}$ is the upper $\alpha/2$ quantile of $t_{n-1}$.
\end{teorema}

\begin{observacion}[Why it matters more for small samples]
Since $t_{n-1,\,\alpha/2} > z_{\alpha/2}$ for every finite $n$ (the tails of $t$ are heavier), the $t$-based interval is always \emph{wider} than the $Z$-based interval one would use if $\sigma$ were known. This difference is substantial for small $n$ (e.g., $t_{9,0.025}\approx 2.262$ vs. $z_{0.025}=1.96$, a difference of over $15\%$) and becomes negligible as $n \to \infty$, since $t_\nu \xrightarrow{d} N(0,1)$.
\end{observacion}

\begin{ejemplo}
 \label{exmp:5.4.1}
A sample of $n=9$ tensile strength measurements of a material yields $\bar{X}=48.5$ MPa and $S=3.2$ MPa. Construct a $95\%$ confidence interval for the mean strength $\mu$, using $t_{8,0.025}\approx 2.306$.
\end{ejemplo}

\begin{solucion}
The standard error is $S/\sqrt{n} = 3.2/3 \approx 1.067$. The interval is
\begin{align}
 \bar{X} \pm t_{8,0.025}\cdot\frac{S}{\sqrt{n}} = 48.5 \pm 2.306\cdot 1.067
 \approx 48.5 \pm 2.46 = [46.04,\ 50.96] \text{ MPa}.
\end{align}
Had we instead (incorrectly) used the Normal quantile $z_{0.025}=1.96$, the interval would have been $48.5 \pm 2.09 = [46.41, 50.59]$: narrower than the actual uncertainty of a sample of only $9$ observations justifies.
\end{solucion}
